\documentclass[11pt,a4paper]{article}

\usepackage[utf8]{inputenc}
\usepackage[T1]{fontenc}
\usepackage{graphicx}
\usepackage{booktabs}
\usepackage{multirow}
\usepackage{amsmath,amssymb}
\usepackage{microtype}
\usepackage[margin=2.5cm]{geometry}
\usepackage[hidelinks]{hyperref}
\usepackage[nameinlink,capitalise,noabbrev]{cleveref}

% TODO: Select target venue template (e.g.\ Elsevier \texttt{elsarticle},
% IEEE \texttt{IEEEtran}) before journal submission.

% Shared notation and acronym definitions for SnatchPhaseBench manuscript.
% Keep symbols consistent across sections; update here when notation changes.

\usepackage{acronym}
\usepackage{xspace}

% --- Acronyms (define with \acrodef for preamble-safe setup) ---
\acrodef{TAS}{temporal action segmentation}
\acrodef{CV}{computer vision}
\acrodef{ML}{machine learning}
\acrodef{IMU}{inertial measurement unit}
\acrodef{MOCAP}{motion capture}
\acrodef{DL}{deep learning}
\acrodef{RNN}{recurrent neural network}
\acrodef{LSTM}{long short-term memory}
\acrodef{GRU}{gated recurrent unit}
\acrodef{TCN}{temporal convolutional network}
\acrodef{ST-GCN}{spatial--temporal graph convolutional network}
\acrodef{IoU}{intersection over union}
\acrodef{IWF}{International Weightlifting Federation}

% --- Mathematical symbols ---
\newcommand{\windowsize}{31}
\newcommand{\featuredim}{99}
\newcommand{\numlandmarks}{33}
\newcommand{\numclasses}{7}
\newcommand{\numvideos}{208}
\newcommand{\numathletes}{70}
\newcommand{\numwindows}{21249}

% --- Placeholder helpers ---
\newcommand{\todosource}[1]{\textit{[TODO: verified citation required---#1]}}
\newcommand{\todofig}[1]{\fbox{\parbox[c][0.22\textheight][c]{0.92\linewidth}{\centering\textit{[Figure placeholder: #1]}}}}
\newcommand{\todotab}[1]{\multicolumn{1}{c}{\textit{[#1]}}}


\title{SnatchPhaseBench: A Reproducible Benchmark for Temporal Phase Segmentation in Olympic Weightlifting from Markerless Pose Estimation}
\author{%
  % TODO: Author names and affiliations after collaboration agreement.
  Anonymous Authors
}
\date{} % omit draft date in submission version

\begin{document}
\maketitle

\begin{abstract}
\noindent
Automatic analysis of Olympic weightlifting technique requires reliable identification of biomechanical phases within each lift attempt.
Expert annotation remains the reference standard but does not scale to large video collections or routine coaching workflows.
Markerless pose estimation offers a practical alternative to instrumented \acs{MOCAP}, yet public benchmarks for temporal phase segmentation in weightlifting remain scarce, and published pipelines are often difficult to reproduce exactly.

This manuscript introduces \emph{SnatchPhaseBench}, an evolving research artifact that couples a curated snatch dataset, a documented preprocessing pipeline, and a reproducibility-oriented evaluation framework.
At the time of writing, we have verified dataset reconstruction (\numwindows{} sliding windows from \numvideos{} videos of \numathletes{} athletes), athlete-level data splits, and a frozen \acs{LSTM} baseline protocol derived from prior thesis work.
Quantitative performance claims in saved source artifacts have \emph{not} yet been re-validated on the original trained checkpoint; benchmark comparisons against alternative temporal models are planned but not reported here.

The intended contribution is therefore methodological and infrastructural: a reproducible benchmark protocol for temporal phase segmentation rather than a single architecture.
Future revisions of this living manuscript will incorporate verified benchmark results, segment-level metrics, and public release details as those experiments are completed.
\end{abstract}

\noindent\textbf{Keywords:} temporal action segmentation; Olympic weightlifting; markerless pose estimation; reproducibility; sports biomechanics; benchmark

\section{Introduction}
\label{sec:introduction}

% CONTENT GUIDANCE:
% - Motivate technical analysis of the snatch for coaching and research.
% - Limitations of manual expert annotation at competition speed.
% - Opportunity of markerless pose for scalable biomechanics.
% - State that this paper introduces SnatchPhaseBench (dataset + protocol + baselines).
%
% TODO: Cite weightlifting biomechanics and sports AI surveys (verify in bibliography.bib).
%
% DO NOT claim novelty or report unverified metrics.

\subsection{Contributions}
% TODO (bullet list, only when verified):
% \begin{itemize}
%   \item Dataset of N videos, M athletes, seven phases.
%   \item Athlete-level evaluation protocol with segment metrics.
%   \item Baseline suite: ...
% \end{itemize}

\subsection{Paper organization}
% TODO: One paragraph outlining Sections~\ref{sec:related}--\ref{sec:conclusion}.

\section{Related Work}
\label{sec:related}

% CONTENT GUIDANCE:
% - Pose estimation in sports (MediaPipe, OpenPose, etc.).
% - Temporal action segmentation (MS-TCN, ASFormer, etc.).
% - Weightlifting / strength sport biomechanics.
% - Gap: limited snatch phase segmentation benchmarks.
%
% TODO: Organize into 3--4 subsections with verified citations.

\subsection{Pose Estimation for Sports}
% TODO

\subsection{Temporal Action Segmentation}
% TODO

\subsection{Weightlifting Biomechanics}
% TODO

\section{Dataset}
\label{sec:dataset}

SnatchPhaseBench comprises annotated monocular snatch attempts with per-frame phase labels, segment interval annotations, and precomputed MediaPipe Pose Landmarker keypoints.
This section summarizes verified dataset statistics and preprocessing outputs; items requiring legal confirmation or additional metadata are marked explicitly.

\subsection{Dataset overview}
\label{sec:dataset:overview}

The dataset contains \numvideos{} snatch attempts performed by \numathletes{} athletes.
Each attempt is represented as a video file path key (raw videos are not bundled in the audited export), frame-wise phase annotations, optional segment intervals, and a CSV of pose keypoints aligned by frame index.
Seven phase classes are used for supervised learning; an additional \texttt{unlabeled} class appears in annotation files but is excluded from training windows when the center frame is unlabeled (\Cref{sec:dataset:preprocessing}).

After sliding-window construction, the processed dataset contains \numwindows{} samples of shape $(\windowsize, \featuredim)$ stored as \texttt{X.npy} with integer labels \texttt{y.npy}.
Independent rebuilding from annotation and keypoint CSV files reproduced the SHA-256 checksums recorded in the source manifest for \texttt{X.npy} and \texttt{y.npy}.

\subsection{Video collection and clip definition}
\label{sec:dataset:videos}

Each record corresponds to a single snatch attempt clip referenced by a relative path \texttt{video\_relpath} (e.g.\ \texttt{athlete\_folder/attempt.mp4}).
The audited repository includes 208 unique \texttt{video\_relpath} entries consistent with 208 keypoint CSV files.

\textsc{TODO (evidence required):} Document competition or training source, recording dates, camera viewpoints, native frame rate, resolution, and barbell visibility criteria.
\textsc{TODO (legal):} Confirm redistribution rights before public release.

\subsection{Phase taxonomy}
\label{sec:dataset:taxonomy}

Phase labels follow a seven-class snatch decomposition used in the source thesis pipeline.
\Cref{tab:phase_taxonomy} lists class identifiers; operational biomechanical definitions (joint-angle or barbell-height rules) were not independently verified during reproduction.

\begin{table}[t]
  \centering
  \caption{Phase taxonomy for supervised learning. Phase~0 is excluded when it appears at the center of a training window.}
  \label{tab:phase_taxonomy}
  \begin{tabular}{clp{0.48\linewidth}}
    \toprule
    ID & Phase name & Notes \\
    \midrule
    0 & \texttt{unlabeled} & Annotation padding / undefined frames; dropped at window center \\
    1 & \texttt{setup} & Pre-lift positioning \\
    2 & \texttt{first\_pull} & Initial barbell separation from platform \\
    3 & \texttt{transition} & Knee--hip coordination change between pulls \\
    4 & \texttt{second\_pull} & Explosive extension phase \\
    5 & \texttt{turnover} & Barbell ascent and arm turnover \\
    6 & \texttt{catch} & Receipt of the barbell overhead \\
    7 & \texttt{recovery} & Stabilization and stand-up \\
    \bottomrule
  \end{tabular}
\end{table}

\textsc{TODO:} Insert biomechanical definitions with citations once confirmed with domain experts; add figure illustrating phases on exemplar lift.

\subsection{Annotation strategy}
\label{sec:dataset:annotations}

Two annotation granularities are provided:
\begin{itemize}
  \item \textbf{Frame labels} (\texttt{master\_frame\_labels.csv}): one row per $(\text{video}, \text{frame})$ with \texttt{phase\_id} and \texttt{phase\_name}.
        The audited export contains 35{,}825 frame rows with no duplicate $(\text{video}, \text{frame})$ keys.
  \item \textbf{Segment labels} (\texttt{master\_segment\_labels.csv} and per-video CSV files): contiguous intervals per phase suitable for segment-level evaluation.
        856 per-video segment files are present in the audited export.
\end{itemize}

Annotation was performed using an interactive tool in the source project (\texttt{annotate\_phases.py} in the student repository).
\textsc{TODO:} Document annotator background, instructions, adjudication protocol, and inter-annotator agreement (not measured in current artifact).

\subsection{Pose representation}
\label{sec:dataset:pose}

Pose is represented using MediaPipe Pose Landmarker with \numlandmarks{} body landmarks.
Each frame stores normalized image coordinates $(x, y, z)$, yielding $\featuredim = 3 \times \numlandmarks$ numeric features per frame.
Visibility scores are available in keypoint CSV headers but are \emph{not} used by the frozen baseline pipeline (\texttt{use\_visibility: false}).

Missing values are handled via interpolation followed by forward and backward filling along each coordinate stream (verified in the frozen preprocessing implementation).
\textsc{TODO:} Report MediaPipe model version (\texttt{pose\_landmarker\_full.task}) and extraction settings once the binary task file is available locally (currently an incomplete export pointer in one snapshot).

\subsection{Train, validation, and test splits}
\label{sec:dataset:splits}

Generalization is evaluated under a fixed \emph{athlete-level} split: all attempts from a given athlete belong exclusively to training, validation, or test.
This prevents near-duplicate lifting styles from leaking across splits through highly correlated windows (stride~1; \Cref{sec:protocol:overlap}).

\Cref{tab:split_stats} summarizes verified split counts from \texttt{athlete\_split.json}.
No athlete or video identifier appears in more than one split (automated validation test in the canonical repository).

\begin{table}[t]
  \centering
  \caption{Verified athlete-level data split. Window counts derive from rebuilt \texttt{meta.csv}.}
  \label{tab:split_stats}
  \begin{tabular}{lrrr}
    \toprule
    Split & Athletes & Videos & Windows \\
    \midrule
    Train & 49 & 145 & 14{,}140 \\
    Validation & 10 & 30 & 3{,}232 \\
    Test & 11 & 33 & 3{,}877 \\
    \midrule
    Total & 70 & 208 & 21{,}249 \\
    \bottomrule
  \end{tabular}
\end{table}

\textsc{TODO:} Report athlete demographics (sex, weight class, performance level) if available and cleared for publication.
\textsc{TODO:} Evaluate multi-fold cross-validation as robustness analysis (planned; not yet executed).

\subsection{Preprocessing and generated tensors}
\label{sec:dataset:preprocessing}

The frozen preprocessing pipeline (\Cref{sec:methods:preprocessing}) constructs sliding windows of length $\windowsize = 31$ with stride~1.
Each window receives the phase label at its center frame; windows whose center is \texttt{unlabeled} are discarded.

Generated artifacts per rebuild:
\begin{itemize}
  \item \texttt{X.npy} $\in \mathbb{R}^{\numwindows \times \windowsize \times \featuredim}$ (\texttt{float32})
  \item \texttt{y.npy} $\in \mathbb{Z}^{\numwindows}$ (\texttt{int64})
  \item \texttt{meta.csv} with columns linking each window to \texttt{video\_relpath}, frame indices, athlete split, and center \texttt{phase\_id}
  \item \texttt{label\_map.csv} mapping phase IDs to names (eight entries including \texttt{unlabeled})
\end{itemize}

\Cref{tab:class_distribution} reports the verified window counts per phase in the rebuilt dataset.
Transition and \texttt{second\_pull} remain the least frequent actionable phases, suggesting class imbalance relevant to metric interpretation.

\begin{table}[t]
  \centering
  \caption{Verified window counts per phase after center-frame labeling (rebuilt dataset).}
  \label{tab:class_distribution}
  \begin{tabular}{lc}
    \toprule
    Phase & Windows \\
    \midrule
    setup & 6{,}017 \\
    first\_pull & 2{,}434 \\
    transition & 673 \\
    second\_pull & 1{,}245 \\
    turnover & 1{,}903 \\
    catch & 1{,}912 \\
    recovery & 7{,}065 \\
    \midrule
    Total & 21{,}249 \\
    \bottomrule
  \end{tabular}
\end{table}

\subsection{Dataset statistics summary}
\label{sec:dataset:stats}

\Cref{tab:dataset_stats} consolidates high-level statistics for quick reference.
Additional columns (frame counts, hours of video, camera types) require metadata not yet verified.

\begin{table}[t]
  \centering
  \caption{Dataset statistics (verified fields only). Extended columns are placeholders pending metadata confirmation.}
  \label{tab:dataset_stats}
  \begin{tabular}{ll}
    \toprule
    Statistic & Value \\
    \midrule
    Athletes & \numathletes{} \\
    Videos (attempts) & \numvideos{} \\
    Annotated frames & 35{,}825 \\
    Training phases (excl.\ unlabeled) & \numclasses{} \\
    Window size / stride & \windowsize{} / 1 \\
    Processed windows & \numwindows{} \\
    Feature dimension per frame & \featuredim{} \\
    Rebuild hash match (\texttt{X.npy}, \texttt{y.npy}) & Verified \\
    Native FPS / resolution & \textit{[TODO: metadata required]} \\
    Public release status & \textit{[Pending legal review]} \\
    \bottomrule
  \end{tabular}
\end{table}

\section{Methods}
\label{sec:methods}

We formulate snatch phase identification as supervised temporal labeling from fixed pose features.
This section documents only the \emph{verified} baseline pipeline; subsections for additional benchmark architectures are reserved with explicit \textsc{TODO} markers.

\subsection{Problem formulation}
\label{sec:methods:problem}

Let each video $v$ yield a pose sequence $\mathbf{X}^{(v)} \in \mathbb{R}^{T \times \featuredim}$ after feature extraction, where $T$ is the number of frames.
Frame-wise ground truth is $y^{(v)}_t \in \{0,\ldots,7\}$ with phase~0 denoting \texttt{unlabeled}.

The baseline consumes sliding windows $\mathbf{w}_i \in \mathbb{R}^{\windowsize \times \featuredim}$ and predicts a single class label $\hat{y}_i$ for window~$i$, supervised by the center-frame label $y_{t(i)}$.
Benchmark extensions will additionally produce frame-wise labels $\hat{y}^{(v)}_t$ and contiguous predicted segments for segment-level metrics (\Cref{sec:protocol:metrics}).

\subsection{End-to-end pipeline}
\label{sec:methods:pipeline}

\Cref{fig:pipeline} summarizes the verified processing chain:
(i)~pose keypoints per frame,
(ii)~alignment with frame labels,
(iii)~sliding-window tensor construction,
(iv)~train-only feature standardization,
(v)~\acs{LSTM} training with athlete-level splits,
(vi)~window-level evaluation on held-out athletes.

\begin{figure}[t]
  \centering
  \todofig{End-to-end pipeline: raw video (optional) $\rightarrow$ MediaPipe keypoints $\rightarrow$ window tensor dataset $\rightarrow$ LSTM classifier $\rightarrow$ window-level metrics. Regenerate from \texttt{docs/figures\_plan.md} F2.}
  \caption{Verified SnatchPhaseBench baseline pipeline. Pose CSV files in the current artifact replace on-the-fly extraction when raw video or the MediaPipe task binary is unavailable.}
  \label{fig:pipeline}
\end{figure}

\subsection{Preprocessing}
\label{sec:methods:preprocessing}

Preprocessing follows the frozen implementation ported to the canonical repository (\texttt{dataset\_builder.py}; see \texttt{docs/FROZEN\_BASELINE.md}).

\paragraph{Alignment.}
Frame labels from \texttt{master\_frame\_labels.csv} are merged with keypoint rows on $(\text{video\_relpath}, \text{frame})$.

\paragraph{Missing pose values.}
Per-coordinate linear interpolation across time is applied, followed by forward and backward filling for remaining gaps.

\paragraph{Window construction.}
For each video, all windows of length $\windowsize = 31$ with stride~1 are extracted.
The supervised label is the phase ID at the center frame (index 15 within the window).
Windows whose center label is \texttt{unlabeled} (ID~0) are removed.

\paragraph{Feature tensor.}
Each window is stored as a float32 array of shape $(\windowsize, \featuredim)$ with $\featuredim = 99$ from $(x,y,z)$ coordinates of \numlandmarks{} landmarks.
Visibility channels are excluded in the baseline configuration.

\paragraph{Standardization.}
Before training, features are standardized using mean and variance computed on \emph{training-split windows only}, then applied to validation and test windows.

\subsection{Baseline model: LSTM classifier}
\label{sec:methods:lstm}

The verified baseline is a single-layer \acs{LSTM} followed by dropout and a linear classification head (\texttt{lstm\_trainer.py} / \texttt{LSTMClassifier}).
Given window input $\mathbf{w} \in \mathbb{R}^{\windowsize \times \featuredim}$, the model outputs logits over \numclasses{} training phases and is trained with cross-entropy loss using class-frequency weights computed on the training split.

\Cref{tab:lstm_hyperparams} lists hyperparameters frozen in \texttt{configs/baseline\_lstm.yaml}.

\begin{table}[t]
  \centering
  \caption{Verified \acs{LSTM} baseline hyperparameters (frozen configuration).}
  \label{tab:lstm_hyperparams}
  \begin{tabular}{ll}
    \toprule
    Setting & Value \\
    \midrule
    Hidden size & 128 \\
    \acs{LSTM} layers & 1 \\
    Dropout & 0.2 \\
    Optimizer & Adam ($\mathrm{lr}=10^{-3}$, weight decay $10^{-4}$) \\
    Batch size & 64 \\
    Max epochs & 40 \\
    Early stopping & Validation macro-F1, patience 8 \\
    Class weighting & Enabled (inverse frequency) \\
    Standardization & Train split only \\
    Random seed & 42 \\
    \bottomrule
  \end{tabular}
\end{table}

The final checkpoint selection criterion monitors validation macro-F1.
This window-level classifier constitutes the initial baseline for SnatchPhaseBench; it is not claimed to be state of the art.

\subsection{Benchmark models (planned)}
\label{sec:methods:benchmark}

The software registry supports plugging additional temporal models without modifying the frozen baseline modules.
The following subsections are placeholders for Phase~3 experiments.
\emph{No architectures below have been trained or evaluated in the canonical benchmark yet.}

\subsubsection{GRU baseline}
\label{sec:methods:gru}
\textsc{TODO:} Describe gated recurrent baseline with matched parameter budget once implemented.

\subsubsection{Temporal convolutional network (TCN)}
\label{sec:methods:tcn}
\textsc{TODO:} Describe dilated causal TCN head producing frame-wise logits; cite verified \acs{TCN} literature when added.

\subsubsection{Multi-stage TCN (MS-TCN)}
\label{sec:methods:mstcn}
\textsc{TODO:} Describe multi-stage refinement architecture for dense segmentation.

\subsubsection{Spatial--temporal graph convolutional network (ST-GCN)}
\label{sec:methods:stgcn}
\textsc{TODO:} Define skeleton graph over MediaPipe landmarks; specify input window and readout for phase labels.

\subsubsection{Transformer encoder}
\label{sec:methods:transformer}
\textsc{TODO:} Specify tokenization of pose frames, positional encoding, and prediction head.

\subsection{Window construction illustration}
\label{sec:methods:windows}

\begin{figure}[t]
  \centering
  \todofig{Sliding window of size 31 with stride 1; center-frame label assignment; unlabeled centers discarded. See \texttt{docs/figures\_plan.md} F3.}
  \caption{Sliding-window labeling protocol (verified). Adjacent windows overlap substantially when stride~$=1$, inducing temporal dependence among training and evaluation samples (\Cref{sec:protocol:overlap}).}
  \label{fig:window_construction}
\end{figure}

\section{Experimental Protocol}
\label{sec:protocol}

% CONTENT GUIDANCE:
% - Athlete-level split (no leakage).
% - Window vs frame vs segment metrics (define mathematically).
% - Overlap / autocorrelation discussion (neutral language).
% - Hardware, seeds, software versions from environment.json.
% - Statistical testing plan for future multi-fold CV.

\subsection{Evaluation Metrics}
% TODO: Reference docs/evaluation_metrics.md equations.
% Window-level (baseline), frame aggregation, segmental F1, edit score.

\subsection{Training Details}
% VERIFIED hyperparameters in configs/baseline_lstm.yaml.

\subsection{Reproducibility}
% TODO: Zenodo DOI, config files, deterministic seeds.

\section{Results}
\label{sec:results}

This section defines the tables and figures required for a complete journal submission.
\emph{No quantitative experimental results are reported at this stage} except dataset statistics already verified in \Cref{sec:dataset}.
Each placeholder explicitly states the experiment that will populate it.

\subsection{Baseline reproduction}
\label{sec:results:baseline}

\Cref{tab:baseline_reproduction} will report checkpoint-based evaluation on the fixed test split (3{,}877 windows, 11 athletes).
Population experiment: load original \texttt{best\_model.pt}, run frozen \texttt{evaluate\_checkpoint} on rebuilt \texttt{X.npy}/\texttt{y.npy}, compare against saved thesis \texttt{classification\_report.json}.

\begin{table}[t]
  \centering
  \caption{Baseline reproduction on held-out test athletes. \textbf{Do not populate until checkpoint validation succeeds.} Rows distinguish thesis artifact values vs.\ reproduced checkpoint evaluation vs.\ optional CPU retrain approximation.}
  \label{tab:baseline_reproduction}
  \begin{tabular}{lccc}
    \toprule
    Metric & Thesis artifact & Checkpoint eval & Retrain (approx.) \\
    \midrule
    Accuracy & \textit{[TODO]} & \textit{[TODO]} & \textit{[TODO]} \\
    Macro precision & \textit{[TODO]} & \textit{[TODO]} & \textit{[TODO]} \\
    Macro recall & \textit{[TODO]} & \textit{[TODO]} & \textit{[TODO]} \\
    Macro F1 & \textit{[TODO]} & \textit{[TODO]} & \textit{[TODO]} \\
    Weighted F1 & \textit{[TODO]} & \textit{[TODO]} & \textit{[TODO]} \\
    Test windows & 3{,}877 & 3{,}877 & 3{,}877 \\
    \bottomrule
  \end{tabular}
\end{table}

\textsc{Experiment gate:} obtain real \texttt{best\_model.pt} binary; confirm \texttt{Matches saved report: YES} in reproduction script output.

\begin{figure}[t]
  \centering
  \todofig{Window-level confusion matrix for LSTM baseline on test split (7$\times$7). Source: regenerated from checkpoint predictions export.}
  \caption{Baseline confusion matrix (placeholder). Populate from \texttt{test\_predictions.csv} after checkpoint validation.}
  \label{fig:confusion_matrix}
\end{figure}

\begin{figure}[t]
  \centering
  \todofig{Training and validation curves (loss, accuracy, macro-F1) for baseline LSTM. Source: \texttt{history.csv} or retrain logs.}
  \caption{Baseline training curves (placeholder).}
  \label{fig:training_curves}
\end{figure}

\subsection{Per-class baseline performance}
\label{sec:results:perclass}

\Cref{tab:baseline_perclass} will report per-phase precision, recall, and F1 on the test split.

\begin{table}[t]
  \centering
  \caption{Per-class test metrics for the \acs{LSTM} baseline. Populate after checkpoint validation.}
  \label{tab:baseline_perclass}
  \begin{tabular}{lccc}
    \toprule
    Phase & Precision & Recall & F1 \\
    \midrule
    setup & \textit{[TODO]} & \textit{[TODO]} & \textit{[TODO]} \\
    first\_pull & \textit{[TODO]} & \textit{[TODO]} & \textit{[TODO]} \\
    transition & \textit{[TODO]} & \textit{[TODO]} & \textit{[TODO]} \\
    second\_pull & \textit{[TODO]} & \textit{[TODO]} & \textit{[TODO]} \\
    turnover & \textit{[TODO]} & \textit{[TODO]} & \textit{[TODO]} \\
    catch & \textit{[TODO]} & \textit{[TODO]} & \textit{[TODO]} \\
    recovery & \textit{[TODO]} & \textit{[TODO]} & \textit{[TODO]} \\
    \bottomrule
  \end{tabular}
\end{table}

\subsection{Benchmark comparison}
\label{sec:results:benchmark}

\Cref{tab:benchmark_comparison} will compare all registered temporal models under the shared protocol (\Cref{sec:protocol:benchmark}).

\begin{table}[t]
  \centering
  \caption{Benchmark model comparison on test athletes. Populate after Phase~3 experiments.}
  \label{tab:benchmark_comparison}
  \resizebox{\linewidth}{!}{%
  \begin{tabular}{lcccccc}
    \toprule
    Model & Window Acc & Macro F1 & Frame Macro F1 & F1@50 & Edit score & Params \\
    \midrule
    LSTM (baseline) & \textit{[TODO]} & \textit{[TODO]} & \textit{[TODO]} & \textit{[TODO]} & \textit{[TODO]} & \textit{[TODO]} \\
    GRU & \textit{[TODO]} & \textit{[TODO]} & \textit{[TODO]} & \textit{[TODO]} & \textit{[TODO]} & \textit{[TODO]} \\
    TCN & \textit{[TODO]} & \textit{[TODO]} & \textit{[TODO]} & \textit{[TODO]} & \textit{[TODO]} & \textit{[TODO]} \\
    MS-TCN & \textit{[TODO]} & \textit{[TODO]} & \textit{[TODO]} & \textit{[TODO]} & \textit{[TODO]} & \textit{[TODO]} \\
    ST-GCN & \textit{[TODO]} & \textit{[TODO]} & \textit{[TODO]} & \textit{[TODO]} & \textit{[TODO]} & \textit{[TODO]} \\
    Transformer & \textit{[TODO]} & \textit{[TODO]} & \textit{[TODO]} & \textit{[TODO]} & \textit{[TODO]} & \textit{[TODO]} \\
    \bottomrule
  \end{tabular}}
\end{table}

\begin{figure}[t]
  \centering
  \todofig{Bar chart or radar plot comparing benchmark models across window and segment metrics. See \texttt{docs/figures\_plan.md} F9.}
  \caption{Benchmark comparison figure (placeholder).}
  \label{fig:benchmark_comparison}
\end{figure}

\subsection{Segment-level evaluation}
\label{sec:results:segment}

\Cref{tab:segment_metrics} will report segmental F1 at IoU thresholds and edit score derived from aggregated frame predictions.

\begin{table}[t]
  \centering
  \caption{Segment-level test metrics. Populate after frame/segment inference pipeline is run for each model.}
  \label{tab:segment_metrics}
  \begin{tabular}{lcccc}
    \toprule
    Model & F1@10 & F1@25 & F1@50 & Edit score \\
    \midrule
    LSTM (baseline) & \textit{[TODO]} & \textit{[TODO]} & \textit{[TODO]} & \textit{[TODO]} \\
    \multicolumn{5}{l}{\textit{[Additional models after benchmark experiments]}} \\
    \bottomrule
  \end{tabular}
\end{table}

\subsection{Ablation studies}
\label{sec:results:ablation}

\Cref{tab:ablation} will summarize ablations defined in \Cref{sec:protocol:ablations}.

\begin{table}[t]
  \centering
  \caption{Ablation study results (placeholder).}
  \label{tab:ablation}
  \begin{tabular}{llcc}
    \toprule
    Factor & Variant & Macro F1 & $\Delta$ vs.\ baseline \\
    \midrule
    Window size & \textit{[TODO]} & \textit{[TODO]} & \textit{[TODO]} \\
    Stride & \textit{[TODO]} & \textit{[TODO]} & \textit{[TODO]} \\
    Visibility features & \textit{[TODO]} & \textit{[TODO]} & \textit{[TODO]} \\
    Standardization & \textit{[TODO]} & \textit{[TODO]} & \textit{[TODO]} \\
    \bottomrule
  \end{tabular}
\end{table}

\subsection{Runtime and efficiency}
\label{sec:results:runtime}

\Cref{tab:runtime} will report training time, inference throughput, and parameter counts.

\begin{table}[t]
  \centering
  \caption{Runtime comparison (placeholder). Requires consistent hardware across runs.}
  \label{tab:runtime}
  \begin{tabular}{lccc}
    \toprule
    Model & Parameters & Train time & Inference ms/window \\
    \midrule
    LSTM (baseline) & \textit{[TODO]} & \textit{[TODO]} & \textit{[TODO]} \\
    \bottomrule
  \end{tabular}
\end{table}

\subsection{Dataset overview figure}
\label{sec:results:dataset_fig}

\begin{figure}[t]
  \centering
  \todofig{Dataset overview: example frames per phase, athlete diversity montage, annotation timeline. See \texttt{docs/figures\_plan.md} F1.}
  \caption{Dataset overview (placeholder). Requires cleared exemplar videos for publication.}
  \label{fig:dataset_overview}
\end{figure}

\begin{figure}[t]
  \centering
  \todofig{Split visualization: athletes assigned to train/val/test without leakage. See \texttt{docs/figures\_plan.md} F4.}
  \caption{Athlete-level split visualization (placeholder).}
  \label{fig:split_visualization}
\end{figure}

\begin{figure}[t]
  \centering
  \todofig{Class distribution bar chart (windows per phase). Can be generated from verified counts in \Cref{tab:class_distribution}.}
  \caption{Window-level class distribution (placeholder; data verified in \Cref{tab:class_distribution}).}
  \label{fig:class_distribution}
\end{figure}

\subsection{Error analysis}
\label{sec:results:error}

\begin{figure}[t]
  \centering
  \todofig{Qualitative error analysis: misclassified transition vs.\ second\_pull examples with pose overlay. See \texttt{docs/figures\_plan.md} error analysis figure.}
  \caption{Qualitative error analysis (placeholder). Populate after baseline predictions are validated.}
  \label{fig:error_analysis}
\end{figure}

\textsc{TODO:} Narrative discussion of systematic errors (transition confusion, pose occlusions) after confusion matrix is populated.

\section{Discussion}
\label{sec:discussion}

This section provides a structured outline for interpreting SnatchPhaseBench results once benchmark experiments are complete.
No performance claims are made at the current manuscript stage.

\subsection{Interpretation of baseline performance}
\label{sec:discussion:interpretation}

\textsc{TODO:} After checkpoint validation, interpret window-level baseline metrics relative to class imbalance (\Cref{tab:class_distribution}) and known difficulty of transition phases.
\textsc{TODO:} Compare checkpoint vs.\ retrain reproduction gap and attribute to hardware, stochasticity, or implementation differences only with evidence.

\subsection{Overlap and metric choice}
\label{sec:discussion:overlap}

\textsc{TODO:} Discuss how stride-1 window evaluation affects apparent accuracy vs.\ segment-level scores.
\textsc{TODO:} Reference temporal autocorrelation analysis in the repository; quantify correlation among adjacent test windows.

\subsection{Benchmark comparison across temporal architectures}
\label{sec:discussion:benchmark}

\textsc{TODO:} Summarize ranking from \Cref{tab:benchmark_comparison} without overclaiming statistical significance until \Cref{sec:protocol:stats} is executed.
\textsc{TODO:} Relate differences to receptive field, parameter count, and suitability for short vs.\ long phase durations.

\subsection{Comparison with prior work}
\label{sec:discussion:prior}

\textsc{TODO:} Position SnatchPhaseBench metrics relative to prior weightlifting studies cited in \Cref{sec:related:wl} once bibliography is complete.
\textsc{TODO:} Clarify that direct numeric comparison may be impossible if splits and labels differ.

\subsection{Practical implications for coaching and sports science}
\label{sec:discussion:practice}

\textsc{TODO:} Discuss whether segment-level accuracy is sufficient for automated feedback tools.
\textsc{TODO:} Address deployment constraints: monocular club footage, real-time latency, interpretability for coaches.

\subsection{Threats to validity}
\label{sec:discussion:validity}

\paragraph{Internal validity.}
\textsc{TODO:} Overlapping windows, single annotation pass, pose estimator bias.

\paragraph{External validity.}
\textsc{TODO:} Athlete cohort diversity, camera conditions, generalization to other weightlifting movements (clean and jerk).

\paragraph{Construct validity.}
\textsc{TODO:} Alignment between operational phase labels and biomechanical theory once definitions are finalized.

\subsection{Future work}
\label{sec:discussion:future}

\textsc{TODO:} Public release, multi-annotator study, paired MOCAP subset, barbell tracking fusion, cross-movement transfer, real-time demo.

\section{Limitations}
\label{sec:limitations}

We report limitations supported by the Phase~1 audit and Phase~2 reproduction effort.
Speculative limitations without evidence are avoided.

\begin{itemize}
  \item \textbf{Single cohort and setting.}
        All verified attempts come from one curated collection (\numvideos{} videos, \numathletes{} athletes) with unknown camera diversity until metadata are confirmed (\Cref{sec:dataset:videos}).

  \item \textbf{Monocular pose inputs.}
        The benchmark relies on 2D/2.5D MediaPipe landmarks rather than multi-view or instrumented kinematics, limiting direct comparison to laboratory \acs{MOCAP} studies.

  \item \textbf{Dependence on pose estimation quality.}
        Errors in occluded or fast phases (turnover, catch) propagate to temporal models; pose error has not been quantified on this dataset.

  \item \textbf{Manual annotation without inter-rater study.}
        Phase labels reflect the source thesis annotation protocol; inter-annotator agreement and adjudication rules are undocumented.

  \item \textbf{Incomplete public artifact in one export.}
        Binary files (\texttt{best\_model.pt}, frozen \texttt{X.npy}/\texttt{y.npy}, MediaPipe task file) were distributed as Git LFS pointers in an exported snapshot, blocking exact checkpoint evaluation until binaries are obtained.

  \item \textbf{Checkpoint validation pending.}
        Thesis metrics preserved in \texttt{classification\_report.json} have not been re-validated on the original trained weights in the canonical reproduction environment.

  \item \textbf{Window-level baseline protocol.}
        Stride-1 sliding windows create strongly correlated samples from the same video; window-level accuracy may overstate independent generalization performance relative to segment-level metrics.

  \item \textbf{Class imbalance.}
        Transition and second pull contain far fewer windows than setup and recovery (\Cref{tab:class_distribution}), affecting macro-averaged metrics.

  \item \textbf{Legal and privacy constraints.}
        Public release requires pseudonymization and clearance of video rights; athlete identifiers in folder names remain sensitive.

  \item \textbf{Single fixed split.}
        One athlete-level partition (49/10/11) may not capture variance in split assignment; cross-validation is planned but not yet executed.

  \item \textbf{Benchmark scope not yet evaluated.}
        Alternative temporal architectures and segment-level benchmark numbers are planned infrastructure only; no comparative scientific conclusions can be drawn yet.
\end{itemize}

\section{Conclusion}
\label{sec:conclusion}

\textsc{TODO:} Summarize verified contributions once benchmark and checkpoint validation are complete.

\paragraph{Verified progress (draft bullet outline).}
\begin{itemize}
  \item SnatchPhaseBench dataset structure, annotations, and athlete-level splits documented and programmatically validated.
  \item Preprocessing pipeline rebuilt processed tensors matching published checksums.
  \item Reproducibility-oriented software infrastructure and evaluation modules implemented.
\end{itemize}

\paragraph{Pending before submission.}
\begin{itemize}
  \item Validate original \acs{LSTM} checkpoint against thesis metrics.
  \item Execute multi-model benchmark experiments and populate \Cref{sec:results}.
  \item Complete literature review with verified citations.
  \item Finalize public dataset release and ethical approvals.
\end{itemize}

\textsc{TODO:} Closing paragraph on expected impact for reproducible sports biomechanics benchmarks—revise after empirical results are available.


\bibliographystyle{plain}
\bibliography{bibliography}

\end{document}
